% ==========================================
% SECTION 1: INTRODUCTION
% ==========================================

\subsection{Motivation and Problem Statement}
\label{sub:motivation}

Public administrations and governmental bodies are under increasing
pressure to operate efficiently, transparently, and cost-effectively
while strictly adhering to regulatory and legal frameworks. However,
traditional public sector processes often suffer from inherent
inefficiencies, such as highly bureaucratic workflows, resource
constraints, and rigid batch-processing habits. These systemic issues
frequently lead to severe process bottlenecks, prolonged throughput
times, and ultimately, a loss of public trust and financial revenue.

A prominent example of such an administrative workflow is the
processing and collection of road traffic fines. When traffic
violations are not processed and collected efficiently, the
administration faces delayed cash flows and increased costs due to
forced legal escalations through credit collection agencies.
Analyzing the well-known \textit{Road Traffic Fine Management}
dataset---extracted from a real-world information system of an Italian
local police force \parencite{deleoni2015}---is therefore highly
relevant. It provides a large-scale, complex event log that closely
reflects the operational reality of public administrations. By
applying advanced process analytics to this dataset, this project
demonstrates how organizations can transition from a reactive,
intuition-based management approach to a proactive, data-driven
optimization strategy.

\subsection{Project Objectives and Research Questions}
\label{sub:objectives}

The primary objective of this project is to comprehensively analyze
the Road Traffic Fine Management process by bridging the gap between
descriptive Process Mining \parencite{vanderaalst2016} and predictive
Machine Learning. Instead of solely discovering how the process was
executed in the past, this report aims to predict future case outcomes
and provide actionable recommendations to actively mitigate
operational risks.

To guide the analytical pipeline, this report addresses the following
core research questions:
\begin{itemize}
    \item \textbf{RQ1 (Process Discovery):} What are the predominant
        control-flow variants, and where do the most critical temporal
        bottlenecks occur within the administration?
    \item \textbf{RQ2 (Operational Behavior):} To what extent does the
        administration exhibit batching behavior, and how do workload
        fluctuations impact the overall process performance?
    \item \textbf{RQ3 (Predictive Monitoring):} Can we accurately
        predict the escalation of a fine to credit collection early in
        the process lifecycle by utilizing Machine Learning algorithms
        such as Random Forest and Long Short-Term Memory networks?
    \item \textbf{RQ4 (Prescriptive Analytics):} How can predictive
        insights and SHAP-based model explainability be leveraged to
        generate prescriptive, rule-based recommendations that
        minimize legal escalation costs?
\end{itemize}

\subsection{Structure of the Report}
\label{sub:structure}

To answer the proposed research questions, this report is structured
sequentially along the data science lifecycle.
Sections~\ref{sec:data_prep} and~\ref{sec:eda} cover data
preparation and exploratory analysis.
Section~\ref{sec:discovery} applies process discovery and conformance
checking to construct and evaluate formal process models.
Section~\ref{sec:predictive} presents the predictive process
monitoring pipeline for both outcome classification and remaining time
prediction. Section~\ref{sec:interpretability} leverages SHAP-based
explainability to derive prescriptive recommendations, while
Section~\ref{sec:generative} describes synthetic event log generation.
Section~\ref{sec:deployment} presents the prototypical Streamlit
dashboard, and Section~\ref{sec:conclusion} concludes with a summary,
limitations, and future work.